Astronomers studied a spiral galaxy with an inclination angle of $90^\circ$
from the plane of the sky (``edge-on'') and apparent magnitude 8.5. They
measured the rotational velocity and radial distance from the Galactic center
and plotted its rotation curve.

\begin{description}
\item[(2.1)] Approximate the rotation curve in Figure 1 with a continuous
  function $V(D)$ composed of two straight lines.
\item[(2.2)] Using the same observations, they estimated that the rotation
  period of the pressure wave in the galactic disk is half of the rotation
  period of the mass of the disk. Estimate the time it takes for a spiral arm
  to take another turn around the galactic center (use the function
  constructed in 2.1).
\item[(2.3)] Calculate the distance to the galaxy using the Tully--Fisher
  relation (see Table of Constants).
\item[(2.4)] Calculate the maximum and minimum values of the observed
  wavelengths of the hydrogen lines corresponding to 656.28\,nm in the
  spectrum of this galaxy. Hint: also take into account the cosmological
  expansion.
\item[(2.5)] Using Figure 1, estimate the mass of the galaxy up to a radius of
  $3 \times 10^3$ light-years.
\item[(2.6)] Estimate the number of stars of the galaxy, assuming that:
  \begin{itemize}
  \item the mean mass of the stars is equal to one solar mass and one third of
    the baryonic mass of the galaxy is in the form of stars, and;
  \item the fraction of baryonic to dark matter in the galaxy is the same as
    the fraction for the whole Universe (see Table of Constants).
  \end{itemize}
\end{description}
